\chapter{Hairy Tongue}

\section*{Definition}
Hairy tongue (lingua villosa) is a benign condition characterized by elongation and discoloration of the filiform papillae on the dorsal tongue surface, creating a hair-like or fur-like appearance with colors ranging from white to yellow, brown, green, or black.

\section*{What Happens in Your Body}
Failure of normal desquamation of the keratinized filiform papillae leads to excessive elongation (up to 15 mm versus normal < 1 mm). Chromogenic bacteria and yeast (Candida, Aspergillus, and pigment-producing actinomycetes) colonize the elongated papillae, producing the characteristic discoloration.

\section*{Causes}
\begin{itemize}
  \item \textbf{Poor oral hygiene:} Inadequate tongue brushing and overall dental care
  \item \textbf{Tobacco use:} Smoking or chewing tobacco contributing to keratin accumulation
  \item \textbf{Xerostomia:} Reduced salivary flow from medications (anticholinergics, antipsychotics) or dehydration
  \item \textbf{Diet:} Soft or liquid diets lacking mechanical abrasive cleaning
  \item \textbf{Antibiotics:} Broad-spectrum antibiotics altering oral flora, allowing chromogen overgrowth
  \item \textbf{Antiseptic mouthwashes:} Prolonged use of oxidizing agents (hydrogen peroxide, chlorhexidine)
  \item \textbf{Radiation therapy:} Head and neck radiation causing xerostomia and altered desquamation
\end{itemize}

\section*{When to See a Doctor}
See your doctor if:
\begin{itemize}
  \item Hairy tongue persists > 4 weeks despite improved oral hygiene
  \item Painful or burning sensation of the tongue
  \item Dysphagia or globus sensation (pharyngeal irritation from elongated papillae)
  \item Halitosis unresponsive to standard oral care
  \item Dysgeusia (altered taste or persistent bad taste)
\end{itemize}

\section*{Related Symptoms}
\begin{itemize}
  \item Elongated filiform papillae on the tongue dorsum
  \item Discoloration (white, yellow, brown, green, black)
  \item Halitosis
  \item Dysgeusia (metallic or bitter taste)
  \item Gagging or tickling sensation on the soft palate
\end{itemize}
\textbf{Important:} Hairy tongue is a purely cosmetic and benign condition with no malignant potential; the color (including black) is due to chromogenic bacteria and not indicative of systemic disease.

\section*{Self-Care / Home Management}
\begin{itemize}
  \item Tongue scraping or brushing twice daily with a soft-bristled toothbrush or tongue scraper
  \item Mechanical debridement: gentle scraping from posterior to anterior, covering the entire dorsum
  \item Increase water intake (6--8 glasses daily) for improved oral hydration
  \item Discontinuation of tobacco products
  \item Chewing sugar-free gum to stimulate salivary flow
  \item Avoid antimicrobial mouthwashes (chlorhexidine, hydrogen peroxide)
  \item Dietary modification: increase fibrous, abrasive foods (raw vegetables, apples)
\end{itemize}

\section*{How Your Doctor Will Evaluate You}
\begin{itemize}
  \item Visual inspection: elongated filiform papillae with discoloration on the dorsal tongue
  \item History taking for medication use (antibiotics, anticholinergics), tobacco, and oral hygiene practices
  \item KOH preparation of scraping to assess for superimposed candidiasis (white form)
  \item Distinguish from oral hairy leukoplakia (OHL): OHL appears on the lateral tongue, non-scrapable, associated with EBV and immunosuppression
  \item Biopsy rarely indicated unless lesion is asymmetric, indurated, or concerning for malignancy
\end{itemize}

\section*{Treatment Options}
\begin{itemize}
  \item Improved oral hygiene and mechanical tongue debridement (scraping/brushing) is first-line and usually effective within 1--2 weeks
  \item Smoking cessation and discontinuation of causative mouthwashes
  \item Topical keratolytic agents: 30--40\% urea solution applied to the tongue for resistant cases
  \item Topical retinoids (tretinoin 0.05\% gel) for refractory hyperkeratosis
  \item Oral antifungal therapy (fluconazole 200 mg daily for 7 days) if concurrent candidiasis confirmed
  \item Reassurance: benign, self-limited condition with excellent prognosis
  \item Referral to dentistry or oral medicine if refractory to conservative measures
\end{itemize}

\section*{References}
\begin{thebibliography}{99}
\bibitem{ref1} Gurvits GE, Tan A. "Black Hairy Tongue Syndrome." World J Gastroenterol. 2014;20(31):10845--10850.
\bibitem{ref2} Reamy BV, Derby R, Bunt CW. "Common Tongue Conditions in Primary Care." Am Fam Physician. 2010;81(5):627--634.
\bibitem{ref3} Neville BW, Damm DD, Allen CM, et al. "Oral and Maxillofacial Pathology." 4th ed. Elsevier; 2016.
\bibitem{ref4} Scully C. "Oral and Maxillofacial Medicine." 3rd ed. Churchill Livingstone, 2013.
\bibitem{ref5} Byrd JA, Bruce AJ, Rogers RS III. "Glossitis and Other Tongue Disorders." Dermatol Clin. 2003;21(1):145--154.
\end{thebibliography}
